\section{Kontekst i cel pracy}

% Oprogramowanie wymaga zmian
% Zmiany są coraz to bardziej kosztowne
% Optymalizacja kosztów jest istotna

Modyfikacja istniejącego kodu źródłowego jest zjawiskiem powszechnym.
Powody takiego stanu rzeczy można podzielić na trzy najistotniejsze elementy.

Pierwszym powodem jest istnienie zwinnego procesów wytwórczych oprogramowania, z
których można wyróżnić metodologię \textit{Scrum} jako najbardziej powszechną
w 2024 roku: 78\% przedsiębiorstw zajmujących się pośrednio lub bezpośrednio
tworzeniem oprogramowania używa właśnie jej~\cite{pwc_agile_report}.
Zakłada ona, że oprogramowanie podlega zmianom ze sprintu (cyklu) na sprint,
a najważniejszym celem jest dostarczanie za każdym razem działającego
oprogramowania.
W związku z tym paradygmat, który zakłada wyłącznie tworzenie nowego kodu
źródłowego do z góry zaprojektowanego systemu, uległ dezaktualizacji.
Przy wprowadzaniu zmian programiści poświęcają więcej czasu na czytanie
i zrozumienie zastanego kodu niż na tworzenie nowego, głównie
ze względu na to, że wstawienie nowej, często nieprzewidzianej wcześniej
funkcjonalności wymaga dostosowania do niej istniejącego kodu źródłowego.

% dehaghani2013factors - 90% kosztów oprogramowania w całym cyklu życia to
% jego utrzymanie (modyfikacja w celu aktualizacji do zmieniających się
% wymagań i otoczenia w kótrym program funkcjonuje).

Drugim powodem jest czas życia oprogramowania po wdrożeniu i dynamicznie
zmieniające się środowisko.
Z badań~\cite{dehaghani2013factors} wynika, że 90\% procent kosztów oprogramowania
należy skategoryzować jako jego utrzymanie, to jest modyfikację funkcjonalności
w obliczu zmieniającego się środowiska, które wymaga nowych integracji z innymi
systemami, zmian otoczenia prawnego, ewoluujących (lub zmienianych za pomocą
szybkiej rewolucji) struktur przedsiębiorstw, i tak dalej.

Trzecim, nie mniej istotnym, powodem powszechności zmian wymagań istniejącego
kodu źródłowego jest rachunek uzyskanej wartości wedle poniesionych kosztów.
Utworzenie nowej funkcjonalności jest często bardziej opłacalne na fundamencie
istniejącego już, wdrożonego i znanego użytkownikom oprogramowania niż
wytworzenie kompletnie nowego programu, nawet jeżeli miałoby realizować ono
wyłącznie nową funkcję.

W obliczu powszechności modyfikacji istniejącego kodu źródłowego nasuwa się
pytanie: czy koszty wprowadzania zmian można optymalizować poprzez odpowiednie
tworzenie oprogramowania w pierwszej kolejności?
Innymi słowy, czy struktura istniejącego programu może mieć wpływ na czas,
który programiści muszą poświęcić na zmianę programu, a tym samym przyspieszyć
wdrożenia modyfikacji?

Biorąc pod uwagę powszechność oprogramowania w każdej z dziedzin ludzkiego
życia, globalną dostępność i relatywnie niski koszt wynikający z wolumenu
danych możliwych do przetworzenia wobec używania pracy ludzi jako alternatywy,
odpowiedź na to pytanie umożliwiałaby organizacjom (nie tylko
przedsiębiorstwom czysto komercyjnym, ale również administracji państwowej,
ośrodkom badawczym i innym) potencjalne wprowadzenie oszczędności, często
wiążące się również z osiągnięciem konkurencyjnej przewagi.

W związku z powyższym, celem niniejszej pracy jest zbadanie zależności pomiędzy
stopniem sprzężenia klas w programie napisanym w języku \texttt{Java}
a czasem wymaganym do wprowadzenia w nim zmiany.
Czas ten przekłada się bezpośrednio na koszt, a to jest kluczowym
aspektem omawianym powyżej.

\section{Struktura pracy}

Złożone, trudne do rozwiązania problemy należy dzielić na mniejsze, które
można relatywnie łatwo rozwiązać.
Tak postąpiono z niniejszą pracą: pytanie badawcze podzielono na etapy, a każdy
z nich stanowi oddzielny rozdział, odpowiednio:

\begin{itemize}
    \item Wstęp --- określa cel pracy, intuicję praktyków, a także przedstawia dotychczasowe badania poczynione w
    obszarze wpływu sprzężenia na koszty utrzymania oprogramowania.
    \item Rozdział 1.\@ --- wprowadza pojęcie sprzężenia i dostępnych metryk, przedstawia ich zastosowanie.
    \item Rozdział 2.\@ --- opracowuje pojęcie tożsamych programów i wskazuje metody modyfikacji sprzężenia
    zachowujące tożsamość programu
    \item Rozdział 3.\@ --- wybiera, wyjaśnia i tworzy projekt metody badawczej.
    \item Rozdział 4.\@ --- przedstawia uzyskane wyniki w toku eksperymentu i poddaje je analizie statystycznej.
    \item Rozdział 5.\@ --- wyciąga wnioski powstałe w wyniku niniejszej pracy.
\end{itemize}

Poniżej rozdziału 5.\@ praca zawiera również bibliografię, wykaz skrótów i symboli,
spisy rysunków, tabel i załączników --- zgodnie z obowiązującymi wymaganiami wobec
prac magisterskich Politechniki Warszawskiej.

\section{Pogląd praktyków na oprogramowanie łatwe w utrzymaniu}
\label{sub:practices}
% Literatura praktyków posiada zasadę "niskie sprzężenie, wysoka spójność"
% Czy zatem sprzężenie klas ma wpływ na koszt zmiany oprogramowania?
% Ktoś to musiał już zbadać, prawda?

Postawione pytania dotyczące kosztów utrzymania oprogramowania nie są niczym nowym.
Praktycy zauważyli ten problem już w latach 80.\@ ubiegłego wieku, wskazując
szereg cech programu łatwego do zrozumienia.
Poniższe elementy są syntezą zaczerpniętą~z późniejszych publikacji mających
zastosowanie do języka Java~\cite{martin2009clean, fowler2009, feathers2004working}.

Najoczywistszą z nich jest odpowiedniość nazw stosowanych w kodzie źródłowym.
Jeżeli program stosuje nazwy zmiennych, klas i metod odpowiednio opisujące
rozwiązywany problem, zrozumienie jego staje się znacznie prostsze, niż gdyby
występowały jedynie nic nie znaczące symbole \texttt{a}, \texttt{b} i \texttt{i}.

Nieco mniej jednoznaczne jest stosowanie odpowiedniej formy dokumentacji.
Tu praktycy wskazują różne metody prezentacji i utrzymywania --- jednocześnie
wskazując, że nadmiar dokumentacji zaciemnia obraz i jest wręcz kontrproduktywny.
Do programistów należy ocena, ile i w jaki sposób tworzona dokumentacja będzie
optymalna dla osób w przyszłości próbujących zrozumieć informacje w niej zawarte.

Najciekawszym (wedle autora tej pracy) i będącym jednocześnie głównym obszarem
zainteresowań niniejszej pracy magisterskiej jest wykorzystanie optymalnej
struktury programu.
Przytoczeni praktycy twierdzą, że utworzenie programu o odpowiedniej strukturze
będzie rzutowało na zrozumienie jego przeznaczenia i sposobu, w jaki modeluje i
obsługuje dziedzinę zadanej funkcjonalności.
We wszystkich wymienionych przypadkach pada sugestia, że programista powinien
utrzymywać moduły (niezależnie, czy są to klasy w przypadku programów obiektowych,
czy procedury w ich strukturalnych odpowiednikach) o wysokiej spójności i niskim
sprzężeniu.

Niniejsza praca ma za zadanie zdefiniować drugi składnik polecenia --- sprzężenie ---
a także zaproponować i przeprowadzić badanie empiryczne, sprawdzając doświadczalnie
zalecenia praktyków.

% Nie? To ja to zbadam.
% Zawartość pracy tak bardzo ogólnie

\section{Dotychczasowe badania}

W obszarze badań nad wpływem stopniem sprzężenia na czas wprowadzania zmian nie
opublikowano zbyt wiele badań spełniających kryteria publikacji naukowej.
Nie odnaleziono też badań eksperymentalnie potwierdzających wpływ sprzężenia
na czas wprowadzania zmian w metodzie zaproponowanej w rozdziale~\ref{cha:metoda}.
Poniżej przytoczone są publikacje podobnych eksperymentów,
które w bezpośredni lub pośredni sposób wpłynęły na zaprojektowanie badania
przeprowadzonego w niniejszej pracy.

W najbardziej odległym, ale wartym wspomnienia, od tej pracy obszarze badawczym
leżą badania polegające na analizie zmian wprowadzanych w projektach otwartego
źródła, stwierdzające o powszechności refaktoryzacji w celu redukcji sprzężenia
pomiędzy klasami~\cite{du2004refactoring}.
Tego typu badania dowodzą, że programiści są zainteresowani zmniejszaniem siły
sprzężenia, ale nie wskazują, czy wpływa to na poniesione koszty utrzymania
oprogramowania.
Pośrednio potwierdzają one przekonania praktyków wspomniane w
podrozdziale~\ref{sub:practices}.

% arisholm2004evaluating - opisuje badanie, w którym sprawdzają strukturę delegowaną wersus scentralizowaną
% i wniosek jest taki, że styl najbardziej zależy od znajomości stylu, bardziej niż od jego samego.
% W badaniu studenci i młodzi programiści lepiej poradzili sobie z wprowadzaniem zmian w stylu scenatralizowanym,
% a zaawansowani w delegującym stylu.
W 2004 roku zostało opublikowane badanie, w którym autorzy przeprowadzają
eksperyment z udziałem studentów, młodszych i starszych programistów.
W ramach eksperymentu wprowadzają oni zmiany do istniejących systemów, które
zostały utworzone w dwóch wariantach: struktury scentralizowanej, która to
jest oparta o jeden moduł wykonujący i moduły stanowiące struktury danych, i
delegowaną, w której każdy moduł ma oddelegowaną odpowiedzialność, przez co
struktura programu jest bardziej deklaratywna.
W ramach eksperymentu mierzono, między innymi, czas wprowadzania zmian.
Wnioskiem z badania jest, że sam styl nie ma większego znaczenia --- młodzi
programiści i studenci lepiej radzili sobie w wariancie scentralizowanym, gdyż
był to sposób tworzenia oprogramowania im bardziej znany, a starsi programiści
szybciej wprowadzali bardziej poprawne zmiany w kodzie o delegowanej
strukturze~\cite{arisholm2004evaluating}.

% deligiannis2004assessing - kod w stylu resposibility-driven i control-oriented (mainframe) napisany w Java,
% dwie grupy studentów (43 w sumie) po trzy zadania i oceniono wzgfledem zrozumienia, effortu i czasu wykonania.
% Nie wykazano statystycznej różnicy pomiędzy czasem wykonania dla obu stylów, wskazano różnice w rozszerzalności
% na korzyść programu w stylu RD
W tym samym roku inne badanie na 43 studentach inżynierii oprogramowania
dowiodło w toku eksperymentu, że kod źródłowy podzielony na klasy wedle
odpowiedzialności wykazuje statystycznie istotną różnicę w łatwości
rozszerzania niż jego alternatywa oparta o delegację kontroli w stylu
\textit{mainframe}.
Programy zostały napisane w języku \texttt{Java}, a studenci zostali
podzieleni na dwie grupy przypisane do wariantów programu i mieli do
wykonania 3 zadania modyfikacji.
Czas wykonania był jednym z metryk eksperymentu, wraz z oceną zrozumiałości
i kosztu wdrożenia~\cite{deligiannis2004assessing}.

Zakrojonym na szeroką skalę eksperymencie w 2011 roku wzięło udział
90 certyfikowanych programistów Java.
Badanie polegało na wykonywaniu zmian w kodzie źródłowym jednego z czterech
wariantów programu o zmodyfikowanych współczynnikach sprzężenia
(\acrshort{CBO}) i spójności (\acrshort{LCOM}): niskim sprzężeniu i niskiej
spójności, niskim sprzężeniu i wysokiej spójności, wysokim sprzężeniu i niskiej
spójności, i wysokim sprzężeniu i wysokiej spójności.
Sprzężenie w wariantach programów zostały przekształcone poprzez zmianę ilości
wywołań metod obcych klas.
Następnie wszystkie warianty poddano ocenie dwóch niezależnych ekspertów pod
kątem weryfikacji, czy programy realizują te same czynności.
Eksperyment potwierdził istotną różnicę w kosztach programu o wysokiej spójności
i niskim sprzężeniu, jednocześnie potwierdzając zależność czasu wprowadzania
zmian od sprzężenia, gdzie różnił się on o 122\% w wariancie o wysokim
sprzężeniu i wysokiej spójności do wariantu o niskim sprzężeniu i wysokiej
spójności~\cite{ramasubbu2011structural}.
Badanie to jest najbliższe eksperymentowi przeprowadzonemu w niniejszej pracy.

% ramasubbu2011structural - badanie spójności (LCOM) i sprzężenia (CBO) klas eksperyment przerprowadzony
% na 90 certyfikowanych programistach Java
% Zadanie wykonywane w parach i indywidualnie
% Tożsamość programów została przeprowadzaona przez dwóch nielzależnych ekspertów.
% Sprzężenie zostało zmodyfikowane za pomocą wywołań metod (brak szczegółów tych modyfikacji)
% Badanie potwierdza hipotezę o mniejszym czasie wprowadzania zmian dla programów o niskim sprzężeniu i wysokiej spójności
% jednocześnie stwierdzając, że nie ma znaczenia praca w parach czy indywidualnie, 122% większy effort jeżeli
% jest wysoka spójność i wysokie sprzężenie w porównaniu do wysokiej spókności i niskiego sprzężenia.

% bavota2013empirical - badanie postrzegania sprzeżenia programu, programiści utożsamiają sprzężenie z zależnością
% znaczeniową pomiędzy klasami w programie obiektowym bardziej niż z wywołaniami pomiędzy programami


\section{Podsumowanie}

We wstępie przedstawiono powody, dla których cel niniejszej pracy ma kluczowe
znaczenie dla wszystkich zajmujących się tworzeniem i utrzymywaniem
oprogramowania.
Przedstawiono strukturę dokumentu, dzieląc pracę na logicznie następujące
etapy badania.
Przytoczono dotychczasową opinię programistów-praktyków, wskazując pewność
stwierdzenia, że aby kod był tani w utrzymaniu, należy utrzymywać klasy
w możliwie najniższym sprzężeniu między nimi i najwyższej spójności w ich wnętrzu.
Przedstawiono publikacje dotyczące badań empirycznych przeprowadzonych w
obszarze sprzężenia i jakości kodu, wskazano interesujące wnioski z nich
płynące, które zostaną wykorzystane w dalszych rozdziałach.